\newpage
\section{模型的建立与求解}
\subsection{问题一的模型的建立和求解}
\subsubsection{具体分析}

问题一要求...，属于评价/分类/预测/优化/机理分析类问题。本问采用... 方法进行求解。

首先对... 进行...，接着利用... 方法...，然后构建... 模型...，最后通过... 得到...。具体分析流程如图\ref{fig:problem1-flow}所示。

\begin{figure}[htbp]
	\centering
	\begin{tabular}{ccc}
		\fbox{\makebox[0.22\textwidth][c]{步骤 1}} &
		\fbox{\makebox[0.22\textwidth][c]{步骤 2}} &
		\fbox{\makebox[0.22\textwidth][c]{步骤 3}} \\[1em]
		\fbox{\makebox[0.22\textwidth][c]{步骤 4}} &
		\fbox{\makebox[0.22\textwidth][c]{步骤 5}} &
		\fbox{\makebox[0.22\textwidth][c]{步骤 6}} \\[1em]
		\fbox{\makebox[0.22\textwidth][c]{步骤 7}} &
		\fbox{\makebox[0.22\textwidth][c]{步骤 8}} &
		\fbox{\makebox[0.22\textwidth][c]{步骤 9}}
	\end{tabular}
	\caption{问题一分析流程图}
	\label{fig:problem1-flow}
\end{figure}

\subsubsection{模型准备}

在进行建模之前，需要...。本节将从数据预处理和算法选择两方面进行准备。

原始数据来源于...，包含... 条记录和... 个字段。经检查发现以下问题：...（缺失值/量纲不统一/异常值/编码问题）。

针对上述问题，按以下步骤进行预处理：第一步，...（如缺失值处理，具体方法 + 参数）。第二步，...（如标准化/归一化，写出公式）。第三步，...（如特征编码/衍生）。

预处理后，数据规模为... 行 $\times$ ... 列，各指标统计量如下表所示。

\begin{table}[htbp]
	\caption{数据预处理前后统计量对比}
	\centering
	\begin{tabular}{ccccc}
		\toprule
		指标 & 预处理前均值 & 预处理前标准差 & 预处理后均值 & 预处理后标准差 \\
		\midrule
		... & ... & ... & ... & ... \\
		\bottomrule
	\end{tabular}
\end{table}

综上所述，数据预处理完成，为后续建模提供了清洁的输入数据。

该问题本质上是一个... 问题，需要在... 条件下...。常用求解算法包括 A 算法、B 算法、C 算法等。本节从维度 1、维度 2、维度 3 三方面进行评估。

\begin{table}[htbp]
	\caption{算法能力对比}
	\centering
	\begin{tabular}{ccccc}
		\toprule
		算法名称 & 维度 1 & 维度 2 & 维度 3 & 稳定性 \\
		\midrule
		算法 A（缩写） & 优 & 良 & 中 & 良 \\
		算法 B（缩写） & 良 & 优 & 良 & 优 \\
		算法 C（缩写） & 中 & 中 & 优 & 中 \\
		\bottomrule
	\end{tabular}
\end{table}

通过对比可以看出，算法 A 在维度 1 方面表现最优，...。因此本文选择算法 A 作为核心求解方法。

综上所述，本节完成了数据预处理和算法选型，接下来将建立数学模型并进行求解。

\subsubsection{模型建立}

针对问题一，本节将从... 角度建立... 模型。具体过程如下。

步骤 1：...

首先，...。设... 为...，则有：
\begin{equation}
	\cdots = \cdots
\end{equation}

其中，... 表示...。

将... 代入...，可得：
\begin{equation}
	\cdots = \cdots
\end{equation}

进一步展开得：
\begin{equation}
	\cdots = \cdots
\end{equation}

步骤 2：...

根据... 的性质，可以建立以下关系：
\begin{equation}
	\cdots = \cdots
\end{equation}

其中，... 为...，... 为...。

由式 (1) 和式 (2) 联立可得：
\begin{equation}
	\cdots = \cdots
\end{equation}

步骤 3：...

步骤 4：... 算法

\begin{table}[htbp]
	\caption{算法参数设置}
	\centering
	\begin{tabular}{cccc}
		\toprule
		参数名 & 参数含义 & 取值 & 选取依据 \\
		\midrule
		... & ... & ... & ... \\
		\bottomrule
	\end{tabular}
\end{table}

综上所述，本节完成了... 模型的建立，得到了...（目标函数/评价体系/预测模型），接下来将基于此进行求解。

\subsubsection{模型求解}

基于上述模型，本节利用... 数据进行求解。求解环境为...，核心参数如表\ref{tab:solution-result}所示。求解得到以下结果：

\begin{table}[htbp]
	\caption{... 结果}
	\label{tab:solution-result}
	\centering
	\begin{tabular}{cccc}
		\toprule
		指标 1 & 指标 2 & 指标 3 & 指标 4 \\
		\midrule
		... & ... & ... & ... \\
		\bottomrule
	\end{tabular}
\end{table}

如图\ref{fig:problem1-flow}所示，...（一句话点出图的结论）。

由表\ref{tab:solution-result}可知，...。从图\ref{fig:problem1-flow}可以看出，...。综合以上结果，...（一段分析总结）。
